\section{合卷：盛世如何制造它自己的边界}

雍正、乾隆时期的国家能力，表现为可以把账册、军机、驻防、漕运和地方官员接入较快的调度链。它让清廷能够经营伊犁、藏区、金川和海疆，也让赈济、河工、盐政和身份管理有了更密集的文书形式。把这些能力称为“盛世”并非全无依据。

但能力的半径从不等于命令的半径。远征依赖道路、牲畜、季节与当地合作者；驻防与设治依赖土地、水源、商路和家庭的再安排；审计和整顿也会把更多征收与责任压到地方。征服的胜利、制度的公布和社会秩序的形成，是三段不能合并的时间。

\subsection{制度得失}

\textbf{所得：} 财政整顿与紧急决策机制提高了调兵、处理亏空和跨区协调的能力，多种边地关系得以纳入同一帝国框架。

\textbf{代价：} 军费、夫役、移民安置和日常供养不断扩大；国家分类有助于管理，也常遮蔽盟旗、寺院、绿洲和家庭的不同经验。

\textbf{留下的压力：} 白莲教战争前夕显出的财政与基层不安，提醒人们扩张并未消除治理成本，只是将其转移到新的社会和地区。
